% This document was generated by an LLM.

\documentclass{essay}

\title{Self-Check Questions:\\ Implicit Differentiation and Differential Equations}
\author{}
\date{}

\begin{document}

\maketitle

\begin{enumerate}

  \item What does the equation \[ F(x,y)=0 \] represent geometrically?

    \textit{Answer.} It represents a \emph{curve} (or more precisely,
    a level set) in the $xy$-plane: the set of all points $(x,y)$ for
    which $F$ vanishes.  For a ``nice'' function $F$ this is
    typically a one-dimensional curve, but it can consist of several
    disconnected pieces, cross itself, or reduce to a single point.

  \item How does an implicit curve differ from an explicit function
    \[ y=f(x)? \]

    \textit{Answer.} An explicit function assigns exactly one
    $y$-value to each $x$; its graph passes the vertical-line test.
    An implicit curve $F(x,y)=0$ imposes only a \emph{relation}
    between $x$ and $y$, so a single value of $x$ may correspond to
    several values of $y$ (e.g.\ both $y=4$ and $y=-4$ lie on the
    circle $x^2+y^2=25$ when $x=3$).

  \item Why can a single implicit equation $F(x,y)=0$ represent
    multiple branches of a curve?

    \textit{Answer.} One equation in two unknowns constrains the
    solution set to a one-dimensional object, but does not force it
    to be a graph.  For a given $x_0$ the equation $F(x_0,y)=0$ can
    have multiple roots, each root corresponding to a different
    branch of the curve stacked above $x_0$.

  \item What assumptions are required for implicit differentiation to be valid?

    \textit{Answer.} Near the point of interest one needs: (i) $F$
    has continuous partial derivatives $F_x$ and $F_y$; (ii) $F_y\neq
    0$ at the point (so the Implicit Function Theorem guarantees a
    local smooth branch $y=f(x)$); (iii) consequently, $y$ is a
    differentiable function of $x$ in a neighbourhood of that point.

  \item How do you compute \[ \frac{dy}{dx} \] from an equation such
    as \[ x^2+y^2=25? \]

    \textit{Answer.} Differentiate both sides with respect to $x$,
    treating $y$ as a function of $x$:
    \[
      2x + 2y\,\frac{dy}{dx} = 0.
    \]
    Solving for $dy/dx$ gives
    \[
      \frac{dy}{dx} = -\frac{x}{y}.
    \]

  \item Express $\frac{dy}{dx}$ in terms of the partial derivatives
    $F_x$ and $F_y$ for a curve defined by $F(x,y)=0$, and derive the
    formula by differentiating both sides with respect to~$x$.

    \textit{Answer.} Differentiating the identity $F(x,y(x))=0$ with
    respect to $x$ via the chain rule:
    \[
      \frac{\partial F}{\partial x} + \frac{\partial F}{\partial
      y}\frac{dy}{dx} = 0.
    \]
    Solving (assuming $F_y\neq 0$):
    \[
      \frac{dy}{dx} = -\frac{F_x}{F_y}.
    \]

  \item Why does differentiating \(y^2\) produce \(2y\,y'\) rather
    than just \(2y\)?

    \textit{Answer.} Because $y$ is itself a function of $x$.  The
    chain rule requires multiplying by the derivative of the inner function:
    \[
      \frac{d}{dx}\bigl[y(x)^2\bigr] = 2y(x)\cdot y'(x).
    \]
    The result $2y$ without $y'$ would only be correct if $y$ were an
    \emph{independent} variable.

  \item What is the geometric interpretation of $\frac{dy}{dx}$
    computed by implicit differentiation, and how does it relate to
    the tangent line to the curve $F(x,y)=0$ at a point?

    \textit{Answer.} $dy/dx$ evaluated at $(x_0,y_0)$ is the
    \emph{slope of the tangent line} to the curve at that point.  The
    tangent line is
    \[
      y - y_0 = \left.\frac{dy}{dx}\right|_{(x_0,y_0)}(x-x_0).
    \]
    It is the best linear approximation to the curve near $(x_0,y_0)$.

  \item At which points does the formula for the implicit derivative
    fail to exist?

    \textit{Answer.} Wherever $F_y = 0$, the formula $dy/dx =
    -F_x/F_y$ is undefined.  Such points are \emph{singular points}
    of the curve; the tangent line may be vertical, the curve may
    cross itself, or the IFT may break down entirely.

  \item How can you determine whether a curve has a vertical tangent
    using implicit differentiation?

    \textit{Answer.} A vertical tangent occurs where $dx/dy = 0$,
    i.e.\ where $-F_y/F_x = 0$, i.e.\ where $F_y = 0$ but $F_x \neq
    0$.  Equivalently, $dy/dx$ is undefined ($F_y=0$) but the curve
    is smooth (a regular point in the $x$-direction).

  \item How do you compute $\frac{d^2y}{dx^2}$ for a curve defined
    implicitly by $F(x,y)=0$? Illustrate the procedure using $x^2+y^2=r^2$.

    \textit{Answer.} Differentiate the expression for $y'$ once more
    with respect to $x$, substituting $y'$ wherever it reappears.

    For $x^2+y^2=r^2$ we found $y'=-x/y$.  Differentiating:
    \[
      y'' = \frac{d}{dx}\!\left(-\frac{x}{y}\right)
      = -\frac{y - x y'}{y^2}
      = -\frac{y - x(-x/y)}{y^2}
      = -\frac{y^2 + x^2}{y^3}
      = -\frac{r^2}{y^3}.
    \]

  \item Given a one-parameter family of curves \[ F(x,y,c)=0, \] why
    does differentiating often eliminate the parameter \(c\)?

    \textit{Answer.} Differentiating with respect to $x$ treats $c$
    as a constant, so its derivative is zero.  The differentiated
    equation involves $x$, $y$, $y'$, and possibly $c$.  One can then
    use the original equation $F=0$ to solve for $c$ and substitute,
    or---if $c$ already drops out after differentiation---the
    parameter is gone directly.

  \item Why does eliminating one arbitrary constant usually produce a
    first-order differential equation?

    \textit{Answer.} One differentiation introduces $y'$ (but no
    higher derivatives) and yields one additional equation.  Using
    the original equation to eliminate the single constant $c$ leaves
    a relation involving only $x$, $y$, and $y'$---precisely a first-order ODE.

  \item Starting from \[ x^2+y^2=c^2, \] derive the corresponding
    differential equation.

    \textit{Answer.} Differentiate with respect to $x$:
    \[
      2x + 2y\,y' = 0 \implies y' = -\frac{x}{y}.
    \]
    The constant $c$ has disappeared; the ODE is $y' = -x/y$.

  \item Starting from \[ y=ce^x, \] derive the corresponding
    differential equation.

    \textit{Answer.} Differentiate:
    \[
      y' = ce^x.
    \]
    But $ce^x = y$, so the ODE is $y' = y$.

  \item Starting from \[ y=x^2+c, \] derive the corresponding
    differential equation.

    \textit{Answer.} Differentiate:
    \[
      y' = 2x.
    \]
    The constant vanishes immediately; the ODE is $y'=2x$.

  \item Given a family of curves depending on $n$ arbitrary
    constants, what systematic procedure eliminates all constants to
    produce a differential equation of order~$n$?

    \textit{Answer.} Write the family as $F(x,y,c_1,\ldots,c_n)=0$.
    Differentiate successively $n$ times with respect to $x$,
    obtaining equations $F^{(1)}=0,\ldots,F^{(n)}=0$ involving
    $y',y'',\ldots,y^{(n)}$ alongside the constants.  The $n+1$
    equations $F=0,F^{(1)}=0,\ldots,F^{(n)}=0$ contain $n$ unknowns
    $c_1,\ldots,c_n$; eliminate them (by substitution or resultants)
    to get a single ODE of order $n$ in $x$, $y$, and $y',\ldots,y^{(n)}$.

  \item Why do arbitrary constants disappear when differentiating?

    \textit{Answer.} The derivative of a constant is zero.  An
    additive constant simply vanishes; a multiplicative constant can
    often be expressed in terms of $y$ (or a lower derivative) using
    the original equation, allowing it to be substituted out.

  \item Is it possible for different curve families to generate the
    same differential equation?

    \textit{Answer.} Yes.  Different families may share the same ODE
    if they are all solution sets of that ODE.  A first-order ODE
    typically has a one-parameter general solution, but it may also
    possess \emph{singular solutions} (e.g.\ envelopes) that do not
    belong to that family yet still satisfy the equation pointwise.

  \item What information is lost when constants are eliminated from a
    solution family?

    \textit{Answer.} The specific member of the family---i.e.\ the
    particular solution.  The constant encodes which initial
    condition was imposed, or equivalently which curve in the family
    the solution lies on.  The ODE retains only the \emph{shape}
    (direction field) shared by all members, not the identity of any
    individual curve.

  \item What is meant by a solution of a differential equation?

    \textit{Answer.} A function $y=\varphi(x)$ defined on some
    interval $I$, differentiable as many times as the order of the
    ODE requires, such that substituting $\varphi$ and its
    derivatives into the equation produces an identity on~$I$.

  \item What is the difference between a particular solution and a
    general solution?

    \textit{Answer.} The \emph{general solution} contains all $n$
    arbitrary constants (for an $n$th-order ODE) and represents the
    entire family of solutions.  A \emph{particular solution} assigns
    specific numerical values to those constants, selecting one
    member of the family---typically the one satisfying given initial
    or boundary conditions.

  \item Why does solving a first-order differential equation
    typically introduce one arbitrary constant?

    \textit{Answer.} Solving requires one integration, which
    introduces one constant of integration.  Equivalently, a
    first-order ODE describes slopes; specifying a slope field
    determines the shape of solution curves up to a vertical shift
    (or other one-parameter adjustment), which is fixed by one
    initial condition.

  \item Why does solving a second-order differential equation
    typically introduce two arbitrary constants?

    \textit{Answer.} Two integrations are needed to pass from $y''$
    back to $y$, each introducing one constant.  Geometrically, a
    second-order ODE prescribes curvature; knowing both position and
    velocity at a point (two conditions) is required to pin down a
    unique solution.

  \item What is the relationship between the order of a differential
    equation and the number of constants in its general solution?

    \textit{Answer.} For a well-behaved ODE, the general solution of
    an $n$th-order equation contains exactly $n$ arbitrary constants.
    This reflects the $n$-fold integration needed to reduce the
    equation to $y$, each step contributing one constant.

  \item What is a separable differential equation? Given \[
    \frac{dy}{dx}=\frac{g(x)}{h(y)}, \] how do you find the general
    solution by separating variables?

    \textit{Answer.} A separable ODE is one that can be written so
    that all $y$-dependence is on one side and all $x$-dependence on
    the other.  For $dy/dx=g(x)/h(y)$, rewrite as
    \[
      h(y)\,dy = g(x)\,dx,
    \]
    then integrate both sides:
    \[
      \int h(y)\,dy = \int g(x)\,dx + C.
    \]
    The result (after solving for $y$ if possible) is the general solution.

  \item Given \[ y'=-\frac{x}{y}, \] recover the family of solution curves.

    \textit{Answer.} Separate variables: $y\,dy = -x\,dx$.  Integrate:
    \[
      \frac{y^2}{2} = -\frac{x^2}{2} + C_1 \implies x^2+y^2 = C \quad
      (C=2C_1>0).
    \]
    The solution family is the set of circles centred at the origin.

  \item How can you verify that a given family of curves satisfies a
    differential equation?

    \textit{Answer.} Substitute $y=f(x,c)$ (and its required
    derivatives, computed explicitly) into the ODE and check that the
    equation holds as an \emph{identity} in $x$ for every admissible
    value of $c$.

  \item What is meant by an integral curve of a differential equation?

    \textit{Answer.} An \emph{integral curve} is a single solution
    curve of the ODE drawn in the $xy$-plane---a curve whose tangent
    direction at every point $(x,y)$ agrees with the slope $f(x,y)$
    prescribed by the equation.  Each particular solution corresponds
    to one integral curve.

  \item What is a slope field, and how does it represent the value of
    $y'$ at each point in the plane?

    \textit{Answer.} A slope field (or direction field) is a plot
    over a grid of points $(x_i,y_j)$ of short line segments, each
    with slope $y'=f(x_i,y_j)$.  It gives a visual picture of all
    directions that integral curves must have; solution curves are
    precisely those that are everywhere tangent to these segments.

  \item Why do all curves in a solution family share the same tangent
    directions at corresponding points?

    \textit{Answer.} Because they all satisfy the same ODE
    $y'=f(x,y)$.  At any point $(x_0,y_0)$, the ODE dictates a unique
    slope $f(x_0,y_0)$, regardless of which member of the family
    passes through that point.  Hence every integral curve through
    $(x_0,y_0)$ must have the same tangent there.

  \item For the equation \[ y'=-\frac{x}{y}, \] what does the slope
    field look like near the coordinate axes?

    \textit{Answer.}
    \begin{itemize}
      \item Near the $y$-axis ($x\approx 0$): slopes are close to
        $0$; segments are nearly horizontal.
      \item Near the $x$-axis ($y\approx 0$): $|{-x/y}|\to\infty$;
        segments become nearly vertical (the formula is undefined on
        the $x$-axis itself).
    \end{itemize}
    This is consistent with the integral curves being circles:
    tangent lines are horizontal at the top/bottom of each circle and
    vertical at its leftmost/rightmost points.

  \item In what sense does a differential equation carry the
    geometric content of an entire curve family, even though it
    contains no free parameter~$c$?

    \textit{Answer.} The ODE encodes the slope at \emph{every} point
    of the plane; any curve that everywhere matches these slopes is
    an integral curve.  Thus the ODE implicitly encodes all members
    of the family simultaneously through the direction field---even
    without naming the constant $c$, it determines which directions
    are allowed, and the family is simply the collection of all
    curves that obey these directions.

  \item Why can a differential equation be viewed as a rule assigning
    a direction to each point in the plane?

    \textit{Answer.} Writing $y'=f(x,y)$, the right-hand side assigns
    to each point $(x,y)$ the number $f(x,y)$, which specifies the
    slope---equivalently, the direction $(1, f(x,y))$ up to
    scaling---of any solution curve passing through that point.  This
    is exactly a \emph{vector field} (or direction field) on the plane.

  \item What is the statement of the Implicit Function Theorem?

    \textit{Answer.}
    Let $F:\R^2\to\R$ be continuously differentiable
    near $(x_0,y_0)$ with $F(x_0,y_0)=0$ and $\frac{\partial
    F}{\partial y}(x_0,y_0)\neq 0$. Then there exist an open interval
    $U\ni x_0$ and a unique continuously differentiable function $f :
    U \to \R$ such that $f(x_0)=y_0$ and \[ F(x,f(x))=0 \quad
    \text{for all} x\in U. \] Moreover, $f'(x)=-F_x(x,f(x))/F_y(x,f(x))$.

  \item How does the Implicit Function Theorem justify implicit differentiation?

    \textit{Answer.} The IFT guarantees that, near a point where
    $F_y\neq 0$, the relation $F(x,y)=0$ does define $y$ as a
    \emph{differentiable} function of $x$.  It is therefore
    legitimate to differentiate the identity $F(x,y(x))=0$
    term-by-term with respect to $x$ using the chain rule, yielding
    $F_x + F_y y'=0$ and hence $y'=-F_x/F_y$.

  \item Under what conditions can an implicit relation $F(x,y)=0$ be
    locally expressed as \[ y=f(x)? \]

    \textit{Answer.} It suffices that $F$ is continuously
    differentiable near the point $(x_0,y_0)$, that $F(x_0,y_0)=0$,
    and that $\frac{\partial F}{\partial y}(x_0,y_0)\neq 0$.  These
    are exactly the hypotheses of the Implicit Function Theorem.

  \item Why can an implicit curve fail to be the graph of a function
    globally but still be a function locally?

    \textit{Answer.} Globally, the zero set of $F$ may wrap around,
    have multiple branches over the same $x$-values, or
    self-intersect---all of which violate the vertical-line test.
    But at any point where $F_y\neq 0$, the IFT cuts out a small
    neighbourhood in which only one branch passes, and on that patch
    the curve \emph{is} the graph of a smooth function.

  \item What happens at points where \[ \frac{\partial F}{\partial y}=0? \]

    \textit{Answer.} The IFT hypothesis fails; we can no longer
    guarantee a local smooth branch $y=f(x)$ near such a point.
    Geometrically, these are \emph{singular points}: the tangent line
    may be vertical (if $F_x\neq 0$), or the curve may cross itself,
    develop a cusp, or be isolated.  The implicit-derivative formula
    $-F_x/F_y$ is undefined there.

  \item Given a two-parameter family \[ F(x,y,c_1,c_2)=0, \] how many
    times must you differentiate to eliminate both constants, and
    what is the order of the resulting differential equation?

    \textit{Answer.} You must differentiate \emph{twice}, producing
    equations $F^{(1)}=0$ and $F^{(2)}=0$ that introduce $y'$ and
    $y''$.  Together with the original equation $F=0$ you have three
    equations for the two unknowns $c_1,c_2$; eliminating them yields
    a second-order ODE.

  \item Starting from \[ y=c_1e^x+c_2e^{-x}, \] derive the
    differential equation satisfied by the family.

    \textit{Answer.} Differentiate twice:
    \[
      y'  = c_1 e^x - c_2 e^{-x}, \qquad
      y'' = c_1 e^x + c_2 e^{-x} = y.
    \]
    Hence $y'' = y$, or $y''-y=0$.

  \item Can every differential equation be obtained by eliminating
    constants from an explicit solution family?

    \textit{Answer.} No.  Some ODEs possess \emph{singular solutions}
    (e.g.\ the envelope of the family) that satisfy the equation but
    cannot be obtained from the general solution by any choice of
    constants.  Furthermore, many ODEs have no closed-form solution
    family at all; they are still well-defined equations, not
    derivable from explicit families.

  \item What does it mean for a differential equation to define a
    family of curves rather than a single curve?

    \textit{Answer.} A first-order ODE specifies the slope $y'$ at
    each point but does not fix which curve passes through a given
    point.  Each choice of integration constant $c$ selects a
    different integral curve; collectively these integral
    curves---one for each $c$---form a one-parameter family that
    ``fills'' the plane (under good conditions without crossings),
    and the ODE governs all of them simultaneously.

  \item How do initial conditions select one member of a solution family?

    \textit{Answer.} An initial condition $y(x_0)=y_0$ requires the
    solution curve to pass through the specific point $(x_0,y_0)$.
    Substituting into the general solution $y=f(x,c)$ gives one
    equation for $c$, which (under Picard--Lindelöf conditions) has a
    unique solution, singling out one member of the family.

  \item Why is \[ y(0)=1 \] sufficient to determine a unique solution
    of many first-order ODEs?

    \textit{Answer.} The general solution of a first-order ODE
    contains exactly one arbitrary constant $c$.  Imposing $y(0)=1$
    yields a single equation $f(0,c)=1$ for $c$.  Provided $f$ is
    smooth in $c$ and the equation has a unique root (which the
    Picard--Lindelöf theorem guarantees locally), $c$ is uniquely
    determined, pinning down a unique solution.

  \item State the Picard--Lindel\"of theorem. What conditions on $f$
    ensure that the initial value problem $y'=f(x,y)$, $y(x_0)=y_0$
    has a unique solution on some interval?

    \textit{Answer.} \textbf{Picard--Lindelöf theorem.}  Suppose
    $f(x,y)$ is continuous on a rectangle $R=\{|x-x_0|\le
    a,\,|y-y_0|\le b\}$ and satisfies a \emph{Lipschitz condition} in
    $y$: there exists $L>0$ such that
    \[
      |f(x,y_1)-f(x,y_2)| \le L|y_1-y_2| \quad \text{for all }
      (x,y_1),(x,y_2)\in R.
    \]
    Then there exists $h>0$ such that the IVP $y'=f(x,y)$,
    $y(x_0)=y_0$ has a \emph{unique} solution on $|x-x_0|\le h$.
    (Continuity of $\partial f/\partial y$ on $R$ is a sufficient
    condition for the Lipschitz hypothesis.)

  \item Why are two conditions typically needed for a second-order ODE?

    \textit{Answer.} The general solution contains two arbitrary
    constants.  Two independent conditions (e.g.\ $y(x_0)=y_0$ and
    $y'(x_0)=v_0$) supply two equations that uniquely determine both
    constants, selecting a unique solution.  Specifying only one
    condition leaves one constant free, yielding infinitely many solutions.

  \item How do initial conditions and boundary conditions both serve
    to determine integration constants, and what distinguishes the two?

    \textit{Answer.} Both impose constraints that fix the free
    constants in the general solution.  The distinction is where the
    conditions are applied:
    \begin{itemize}
      \item \emph{Initial conditions} specify the values of $y$ (and
        its derivatives up to order $n-1$) all at the \emph{same}
        point $x_0$.  They match the physical setup of knowing the
        state of a system at one instant, and existence/uniqueness is
        guaranteed by Picard--Lindelöf.
      \item \emph{Boundary conditions} specify values of $y$ (or its
        derivatives) at \emph{different} points---typically the two
        endpoints of an interval.  They arise naturally in spatial
        problems; existence and uniqueness need not hold in general
        and require separate analysis (Sturm--Liouville theory, etc.).
    \end{itemize}

  \item Can you explain, without calculations, why implicit
    differentiation and solving differential equations are inverse processes?

    \textit{Answer.} Starting from a family of curves $F(x,y,c)=0$,
    \emph{implicit differentiation} eliminates the constant $c$ and
    condenses the entire family into a single ODE---moving from many
    curves to one equation.  Conversely, \emph{solving} the ODE
    reintroduces a constant of integration and expands the single
    equation back into a family of curves---moving from one equation
    to many curves.  Differentiation destroys information (which
    specific curve) and produces the ODE; integration restores that
    information (up to an arbitrary constant) and produces the
    family.  The two operations are thus inverse to each other in the
    sense that ``differentiate-then-integrate'' recovers the family
    (up to relabelling of $c$), and ``integrate-then-differentiate''
    recovers the ODE.

\end{enumerate}

\end{document}
